%!TEX TS-program = xelatex
%!TEX encoding = UTF-8 Unicode
% Awesome CV LaTeX Template for Cover Letter
%
% This template has been downloaded from:
% https://github.com/posquit0/Awesome-CV
%
% Authors:
% Claud D. Park <posquit0.bj@gmail.com>
% Lars Richter <mail@ayeks.de>
%
% Template license:
% CC BY-SA 4.0 (https://creativecommons.org/licenses/by-sa-4.0/)
%


%-------------------------------------------------------------------------------
% CONFIGURATIONS
%-------------------------------------------------------------------------------
% A4 paper size by default, use 'letterpaper' for US letter
\documentclass[11pt, a4paper]{awesome-cv}

% Configure page margins with geometry
\geometry{left=1.4cm, top=.8cm, right=1.4cm, bottom=1.8cm, footskip=.5cm}

% Color for highlights
\colorlet{awesome}{awesome-lavender}

% Set false if you don't want to highlight section with awesome color
\setbool{acvSectionColorHighlight}{true}

% If you would like to change the social information separator from a pipe (|) to something else
\renewcommand{\acvHeaderSocialSep}{\quad\textbar\quad}

% Increase letter title font size
\renewcommand*{\lettertitlestyle}[1]{{\fontsize{12pt}{1.2em}\bodyfontlight\bfseries\color{darktext} \underline{#1}}}

% Override letter title: swap application line and recipient, remove date/address
% Override letter closing: same font for "Sincerely," and name, single line break
\makeatletter
\renewcommand{\makelettertitle}{%
  \vspace{8.4mm}
  \setlength\tabcolsep{0pt}
  \setlength{\extrarowheight}{0pt}
  \begin{tabular*}{\textwidth}{@{\extracolsep{\fill}} L{\textwidth - 4.5cm} R{4.5cm}}
    \lettertitlestyle{\@lettertitle} & \letterdatestyle{\@letterdate}
  \end{tabular*}
  \recipienttitlestyle{\@recipientname}\\
  \recipientaddressstyle{\@recipientaddress} \\\\
  \lettertextstyle{\@letteropening}
}
\renewcommand{\makeletterclosing}{%
  \vspace{3.4mm}
  \lettertextstyle{\@letterclosing}\\
  \lettertextstyle{\@firstname\ \@lastname}%
  \ifthenelse{\isundefined{\@letterenclosure}}
    {\\}
    {%
      \\\\\\
      \letterenclosurestyle{\@letterenclname: \@letterenclosure} \\
    }%
}
\makeatother


%-------------------------------------------------------------------------------
%	PERSONAL INFORMATION
%-------------------------------------------------------------------------------
\name{Spas}{Zahariev}
\position{Software Engineer Lead}
\address{Zurich, Switzerland}

\mobile{(+41) 76-212-0497}
\email{spas.zah@gmail.com}
\homepage{https://spaszah.site}
\github{SpasZahariev}
\linkedin{spas-zahariev}


%-------------------------------------------------------------------------------
%	LETTER INFORMATION
%-------------------------------------------------------------------------------
% The company being applied to
\recipient
  {}
  {}
% The date on the letter, default is the date of compilation
\letterdate{}
% The title of the letter
\lettertitle{Application for Software Engineer III, AI/ML, DSP/HTP Optimisation}
% How the letter is opened
\letteropening{Dear Google Recruitment Team,}
% How the letter is closed
\letterclosing{Sincerely,}
% Any enclosures with the letter
\letterenclosure[Attached]{Curriculum Vitae}


%-------------------------------------------------------------------------------
\begin{document}

% Print the header with above personal information
\makecvheader[R]

% Print the footer with 3 arguments(<left>, <center>, <right>)
\makecvfooter
  {\today}
  {Spas Zahariev~~~\textperiodcentered~~~Cover Letter}
  {}

% Print the title with above letter information
\makelettertitle

%-------------------------------------------------------------------------------
%	LETTER CONTENT
%-------------------------------------------------------------------------------
\begin{cvletter}

\lettersection{Why Me?}
I am applying for the Software Engineer III, AI/ML, DSP/HTP Optimisation role with 7+ years of production engineering experience, hands-on AI/ML delivery in Python, and a record of turning performance bottlenecks into measured improvements. At JPMorgan Chase, I cut messaging turnaround by approximately 85\%, from 30 minutes to under 5 minutes at 150K+ messages per hour, by redesigning a Java service around Kafka and asynchronous processing. That is the engineering pattern I would bring to this role: profile the critical path, understand the system constraints, make targeted changes, and validate the result at realistic scale.

At EPAM, I built a Python RAG agent that saves investment committee members an average of 2 hours of preparation per meeting, and I lead a seven-engineer team responsible for a cloud data pipeline that reconciles 500K+ entities per day across 12 source databases. My Google Professional ML Engineer certification complements practical experience with model deployment, evaluation, data processing, MLOps, Vertex AI, BigQuery, TensorFlow/Keras, Docker, and Kubernetes. Outside client work, I customize Linux environments and run local models with llama.cpp, making compute, memory, and hardware constraints tangible rather than abstract.

\lettersection{Why Google?}
The Platforms and Devices challenge is compelling because model quality alone is not enough: useful on-device AI must also meet strict latency, power, thermal, and stability limits. I want to work where profiling data drives hardware-specific optimization, and where software decisions can shape efficiency standards for HTP, DSP, and next-generation devices.

I also bring the versatility and leadership expected of a Google engineer. I have owned architecture through production, worked across backend, cloud, data, and ML systems, and mentored engineers while remaining hands-on. I would contribute a disciplined optimization mindset, production Python and ML infrastructure experience, and the systems breadth needed to help make advanced AI features efficient and dependable on real devices.

\end{cvletter}


%-------------------------------------------------------------------------------
% Print the signature and enclosures with above letter information
\makeletterclosing

\end{document}
